\begin{textsample}{NEG Dim 6 – vem\_ed\_16 – Score -8.00 – vem\_ed\_16/t070.txt}  \label{ex:f6_neg_015}
VEm Virtual Exchange Medium Como vê as iniciativas de cooperação Sul-Sul nos Intercâmbios Virtuais?
É explícita, nas instituições latino-americanas, por exemplo, a resistência a cooperação com o Sul Global.
Uma breve análise do número de acordos de nossas IES com cada região do planeta poderá, facilmente, mostrar que os países do Norte, com suas instituições centenárias e suas ofertas tentadoras, estão na dianteira do interesse de nossos estudantes e pesquisadores como destino de intercâmbio.
A \textbf{possibilidade} dos Intercâmbios Virtuais entre instituições do Sul abre, pois, uma via a ser explorada para uma quebra inicial – que poderá ser seguida de um desejável equilíbrio – da primazia do Norte em nossas \textbf{relações} bilaterais.
As trocas virtuais permitirão o desenvolvimento de algo que é fundamental para a cooperação Sul-Sul, a saber, o conhecimento mútuo dos parceiros \textbf{potenciais}, de sua \textbf{atuação} na pesquisa, de sua cultura.
Somente esse conhecimento permitirá a construção de sólidas parcerias entre instituições do Sul, talvez com uma \textbf{reciprocidade} maior que aquela que temos observado nas \textbf{relações} com o Norte.
Comente a conquista do selo FAUBAI-BRaVE pelo CPS, em outubro de 2022.
O selo FAUBAI-BRaVE foi criado para atestar a participação das IES membros da FAUBAI nas nossas promoções de Intercâmbio Virtual nos moldes do Programa FAUBAI-BRaVE.
A atribuição do selo ao Centro Paula Souza, além de reconhecer a qualidade de seus projetos de Intercâmbio Virtual, representa também o reconhecimento da FAUBAI pela participação do CPS, juntamente com a UNESP e a UFPE, na experiência piloto que construiu esse programa.
Como as demais IES portadoras desse selo, o CPS poderá utilizá-lo para comprovar sua participação no programa diante dos parceiros que vier a conquistar e, também, terá direito a condições especiais nas promoções que o programa FAUBAI-BRaVE vier a lançar no futuro, como ofertas de parcerias estabelecidas pela FAUBAI, gratuidades ou descontos em formações específicas para a área, entre outros.
Como vislumbra o futuro da internacionalização?
O cenário macropolítico internacional aparenta um recrudescimento de ideias totalitárias e, por vezes, nacionalistas, contra o ideal democrático que faz transpor fronteiras.
Essa impressão acentuou-se no período da pandemia, quando as trocas, num primeiro momento, pareciam ser impossíveis.
Entretanto, podemos dizer sem exagero que o braço tecnológico da globalização se estabeleceu solidamente em todas as culturas, permitindo-nos, inclusive, enfrentar com menos sofrimento as agruras das quarentenas.
Espantosamente, as tecnologias que, pelo alto custo do acesso à Web e a equipamentos, sempre foram excludentes, tornaram-se um fator de inclusão de primeira \textbf{ordem} no mundo da educação, para nos restringirmos ao nosso campo.
Isso mostra que detemos um conhecimento \textbf{capaz} de nos fazer superar qualquer obstáculo, incluindo o negacionismo e o obscurantismo que compõem as receitas totalitárias.
Vejo, pois, o futuro da educação superior internacionalizada sempre mais adepto das trocas virtuais, que já provaram ser de grande utilidade, mesmo diante do mau uso que pode, por vezes, permeá-la.
Ter o mundo em nossa casa, com o recurso quase mágico das novas tecnologias, é uma oportunidade que não pode ser \textbf{ignorada}, pois, além dos contatos rápidos e \textbf{constantes}, oferece também uma via mais segura e econômica dada a necessidade de redução da emissão de carbono.
Não se trata de imaginar o futuro da internacionalização apenas em um metaverso \textbf{crescente}, mas sim de integrar essas \textbf{possibilidades} virtuais ao processo educacional e à produção científica.

% matched lemmas: atuação, capaz, constante, crescente, ignorar, ordem, possibilidade, potencial, reciprocidade, relação
\end{textsample}
